\documentclass[11pt]{article}
\usepackage{amsmath,amssymb,amsthm}
\usepackage[margin=1in]{geometry}
\newtheorem{theorem}{Theorem}
\newtheorem{lemma}[theorem]{Lemma}
\title{Problem 365: A Huge Binomial Coefficient}
\author{Project Euler Solutions}
\date{}
\begin{document}
\maketitle

\section{Problem Statement}

Let $n = 10^{18}$ and $k = 10^{9}$. For each triple of distinct primes $(p_1, p_2, p_3)$ with $1000 < p_i < 5000$, compute $\binom{n}{k} \bmod (p_1 p_2 p_3)$ via the Chinese Remainder Theorem. Find the sum of all such results over all $\binom{574}{3}$ triples.

\section{Mathematical Foundation}

\begin{theorem}[Lucas' Theorem]
Let $p$ be prime. Write $n = \sum_{i} n_i p^i$ and $k = \sum_{i} k_i p^i$ in base~$p$. Then
\[
\binom{n}{k} \equiv \prod_{i} \binom{n_i}{k_i} \pmod{p},
\]
where $\binom{n_i}{k_i} = 0$ if $k_i > n_i$.
\end{theorem}

\begin{proof}
In $\mathbb{F}_p[x]$, $(1+x)^p = 1 + x^p$ since $\binom{p}{j} \equiv 0 \pmod{p}$ for $1 \le j \le p-1$. By induction, $(1+x)^{p^i} = 1 + x^{p^i}$. Therefore
\[
(1+x)^n = \prod_i \bigl(1 + x^{p^i}\bigr)^{n_i} \pmod{p}.
\]
Extracting the coefficient of $x^k = x^{\sum k_i p^i}$ from both sides gives $\binom{n}{k} \equiv \prod_i \binom{n_i}{k_i}$.
\end{proof}

\begin{theorem}[Chinese Remainder Theorem]
Let $m_1, m_2, m_3$ be pairwise coprime. Set $M = m_1 m_2 m_3$ and $M_i = M/m_i$. Given residues $r_i$, the unique solution modulo~$M$ is
\[
x \equiv \sum_{i=1}^{3} r_i \cdot M_i \cdot (M_i^{-1} \bmod m_i) \pmod{M}.
\]
\end{theorem}

\begin{proof}
Since $\gcd(M_i, m_i) = 1$, the inverse $M_i^{-1} \bmod m_i$ exists. The term $r_i M_i (M_i^{-1} \bmod m_i)$ is congruent to $r_i$ modulo~$m_i$ and to~$0$ modulo~$m_j$ for $j \neq i$. Summing gives the solution. Uniqueness follows since $|\mathbb{Z}/M\mathbb{Z}| = \prod |\mathbb{Z}/m_i\mathbb{Z}|$.
\end{proof}

\begin{lemma}[Prime Count]
There are exactly $574$ primes in $(1000, 5000)$, giving $\binom{574}{3} = 31{,}436{,}024$ triples.
\end{lemma}

\section{Algorithm}

\begin{verbatim}
function solve():
    primes = sieve(1000, 5000)   // 574 primes
    for p in primes:
        r[p] = lucas(10^18, 10^9, p)
    S = 0
    for each triple (p1, p2, p3):
        S += CRT(r[p1], r[p2], r[p3], p1, p2, p3)
    return S
\end{verbatim}

\section{Complexity Analysis}

\begin{itemize}
\item \textbf{Time:} $O(\pi^3)$ where $\pi = 574$. Each CRT evaluation is $O(1)$; total $\approx 3.1 \times 10^7$ operations.
\item \textbf{Space:} $O(\pi)$ for the residue array.
\end{itemize}

\section{Answer}

\[
\boxed{162619462356610313}
\]

\end{document}
